\documentclass[aspectratio=169,10pt]{beamer}
\usetheme[numbering=fraction,progressbar=frametitle]{metropolis}
\usepackage{fontspec}
\setsansfont{Latin Modern Sans}
\usepackage{amsmath,amssymb,mathtools,graphicx,booktabs}
\usepackage{polyglossia}
\setdefaultlanguage{french}
\setbeamertemplate{navigation symbols}{}
\setbeamerfont{frametitle}{size=\large}
\setbeamerfont{normal text}{size=\normalsize}
\setlength{\parskip}{.45em}
\hypersetup{pdfauthor={},pdfsubject={Lecture slides},pdfkeywords={},pdftitle={Croisement de deux variables qualitatives}}
\providecommand{\N}{\mathbb{N}}
\providecommand{\Z}{\mathbb{Z}}
\providecommand{\D}{\mathbb{D}}
\providecommand{\Q}{\mathbb{Q}}
\providecommand{\R}{\mathbb{R}}
\providecommand{\C}{\mathbb{C}}
\providecommand{\vect}[1]{\overrightarrow{#1}}
\providecommand{\vi}{\vec{\imath}}
\providecommand{\vj}{\vec{\jmath}}
\providecommand{\vk}{\vec{k}}
\providecommand{\norm}[1]{\left\lVert #1\right\rVert}
\providecommand{\vc}[2]{\begin{pmatrix}#1\\#2\end{pmatrix}}
\providecommand{\deter}[2]{\det\left(#1,#2\right)}
\providecommand{\abs}[1]{\left\lvert #1\right\rvert}
\providecommand{\dd}{\,\mathrm{d}}
\providecommand{\e}{\mathrm{e}}
\providecommand{\Card}{\operatorname{Card}}
\providecommand{\E}{\operatorname{E}}
\providecommand{\V}{\operatorname{V}}
\providecommand{\degre}{^{\circ}}
\providecommand{\un}[1]{\,\mathrm{#1}}
\providecommand{\sci}[2]{#1\times10^{#2}}
\providecommand{\pgcd}{\operatorname{PGCD}}
\providecommand{\ppcm}{\operatorname{PPCM}}
\title{Croisement de deux variables qualitatives}
\subtitle{Cours et exemples}
\date{}
\author{}
\begin{document}
\begin{frame}[plain]
\vfill
{\usebeamerfont{title}\inserttitle\par}
\medskip
{\usebeamerfont{subtitle}\insertsubtitle\par}
\medskip
{\usebeamercolor[fg]{progress bar}\rule{\textwidth}{.5pt}}
\vfill
\end{frame}
\begin{frame}{Tableaux croisés et fréquences marginales}
Une variable qualitative décrit une catégorie. Elle est nominale si ses modalités n'ont pas d'ordre naturel (moyen de transport), ordinale si elles sont ordonnées (degré de satisfaction).
Un tableau croisé compte les individus selon deux variables. Les sommes des lignes et colonnes sont les marges. Une fréquence marginale divise une marge par l'effectif total.
\end{frame}
\begin{frame}{Tableaux croisés et fréquences marginales}
Dans une population de \(100\) individus, on observe deux caractères \(A,B\) :
\end{frame}
\begin{frame}{Correction : Tableaux croisés et fréquences marginales}
\(f(A)=32/100=0{,}32\) ; \(f(B)=35/100=0{,}35\) ; \(f(\overline B)=65/100=0{,}65\). La fréquence de \(A\cap B\) est \(20/100=0{,}20\) ; ce n'est pas une fréquence marginale.
\end{frame}
\begin{frame}{Tableaux croisés et fréquences marginales}
Une classe de Première de \(32\) élèves est répartie ainsi. Compléter les marges et calculer les fréquences marginales « NSI » et « filles ».
\end{frame}
\begin{frame}{Correction : Tableaux croisés et fréquences marginales}
\(f(\mathrm{NSI})=21/32=65{,}625\,\%\) ; \(f(\mathrm{filles})=14/32=43{,}75\,\%\).
\end{frame}
\begin{frame}{Fréquences conditionnelles et lecture de phrases}
La fréquence de \(A\) parmi les individus de \(B\), si cette sous-population est non vide, est
\[f_B(A)=\dfrac{\text{effectif vérifiant }A\text{ et }B}{\text{effectif de }B}.\]
Le dénominateur correspond à la population introduite par « parmi ».
\end{frame}
\begin{frame}{Fréquences conditionnelles et lecture de phrases}
Dans le tableau NSI, calculer la fréquence de filles parmi les élèves sans NSI, puis la fréquence des élèves sans NSI parmi les filles. Les deux résultats sont-ils égaux ?
\end{frame}
\begin{frame}{Correction : Fréquences conditionnelles et lecture de phrases}
Parmi les \(11\) élèves sans NSI, \(5\) sont des filles : \(5/11\approx45{,}45\,\%\). Parmi les \(14\) filles, \(5\) ne font pas NSI : \(5/14\approx35{,}71\,\%\). Le numérateur est le même, les populations de référence diffèrent.
\end{frame}
\begin{frame}{Fréquences conditionnelles et lecture de phrases}
Dans le tableau de \(100\) individus, calculer \(f_B(A)\), \(f_A(B)\) et \(f_{\overline A}(B)\). Traduire chaque résultat par une phrase.
\end{frame}
\begin{frame}{Correction : Fréquences conditionnelles et lecture de phrases}
\(f_B(A)=20/35=4/7\approx57{,}14\,\%\) des individus de \(B\) vérifient \(A\). \(f_A(B)=20/32=62{,}5\,\%\) des individus de \(A\) vérifient \(B\). \(f_{\overline A}(B)=15/68\approx22{,}06\,\%\) des individus ne vérifiant pas \(A\) vérifient \(B\).
\end{frame}
\begin{frame}{Fréquences conditionnelles et lecture de phrases}
« \(60\,\%\) des cyclistes sont des élèves » permet-il de conclure que \(60\,\%\) des élèves sont cyclistes ?
\end{frame}
\begin{frame}{Correction : Fréquences conditionnelles et lecture de phrases}
Non. Par exemple, sur \(10\) cyclistes, \(6\) sont élèves ; si la population compte \(100\) élèves au total, seulement \(6\,\%\) des élèves sont cyclistes.
\end{frame}
\begin{frame}{Compléter un tableau à partir de fréquences}
Un groupe compte \(200\) personnes, dont \(120\) utilisent un bus (\(B\)). Parmi ces \(120\) personnes, \(75\,\%\) possèdent un abonnement (\(A\)). Parmi les \(80\) autres, \(25\,\%\) possèdent un abonnement.

\begin{enumerate}
\setcounter{enumi}{0}\item Construire le tableau croisé \(B/\overline B\) et \(A/\overline A\) avec ses marges.
\item Calculer la fréquence marginale de \(A\).
\item Calculer la fréquence de \(B\) parmi \(A\).

\end{enumerate}
\end{frame}
\begin{frame}{Correction : Compléter un tableau à partir de fréquences}
Dans \(B\), \(120\times0{,}75=90\) personnes vérifient \(A\), donc \(30\) ne le vérifient pas. Dans \(\overline B\), \(80\times0{,}25=20\) vérifient \(A\), et \(60\) ne le vérifient pas.
centertabularc|rr|r\&\(A\)\&\(\overline A\)\&Total

 \(B\)\&90\&30\&120

\(\overline B\)\&20\&60\&80

 Total\&110\&90\&200tabularcenter
\(f(A)=110/200=55\,\%\). \(f_A(B)=90/110=9/11\approx81{,}82\,\%\).
\end{frame}
\begin{frame}{Compléter un tableau à partir de fréquences}
Un titre annonce « Trois abonnés sur quatre prennent le bus », en s'appuyant sur \(f_B(A)=75\,\%\). Expliquer l'erreur et proposer une phrase correcte.
\end{frame}
\begin{frame}{Correction : Compléter un tableau à partir de fréquences}
Le titre inverse le conditionnement. La phrase correcte est : « Trois personnes sur quatre parmi celles qui prennent le bus possèdent un abonnement. » La proportion des abonnés qui prennent le bus vaut ici environ \(81{,}82\,\%\).
\end{frame}
\begin{frame}[fragile]{Python}
\small
\begin{verbatim}
# Exemple fictif : (transport, satisfaction)
donnees = [("bus", "oui"), ("marche", "oui"),
           ("bus", "non"), ("velo", "oui"),
           ("bus", "oui"), ("marche", "non")]

bus = [r for r in donnees if r[0] == "bus"]
bus_satisfaits = [r for r in donnees
                 if r[0] == "bus" and r[1] == "oui"]
autres = [r for r in donnees if not r[0] == "bus"]
\end{verbatim}
\end{frame}
\begin{frame}{Traiter un fichier et recueillir des données réelles}
Donner les effectifs des trois listes. Calculer la fréquence de satisfaction parmi les usagers du bus. Remplacer \texttt{and} par \texttt{or} dans le second filtre : combien de lignes sont alors sélectionnées ?
\end{frame}
\begin{frame}{Correction : Traiter un fichier et recueillir des données réelles}
Effectifs : \(3,2,3\). Parmi les usagers du bus, la fréquence de satisfaction vaut \(2/3\). Avec \texttt{or}, on sélectionne les \(3\) usagers du bus et les \(2\) personnes satisfaites utilisant un autre moyen : \(5\) lignes, sans compter deux fois les lignes qui vérifient les deux critères.
\end{frame}
\begin{frame}{Traiter un fichier et recueillir des données réelles}
Recueillir anonymement deux variables : moyen de transport principal et durée du trajet par catégories ordonnées (\(<15\) min ; de \(15\) à moins de \(30\) min ; \(\geq30\) min). Une ligne par réponse, sans nom. Construire le tableau croisé, puis deux fréquences conditionnelles de sens opposés. Présenter effectif, date de collecte et éventuelles non-réponses.
\end{frame}
\begin{frame}{Traiter un fichier et recueillir des données réelles}
Réaliser la collecte et rédiger deux conclusions justifiées. Expliquer pourquoi les résultats ne décrivent pas automatiquement tous les élèves de Madagascar.
\end{frame}
\begin{frame}{Correction : Traiter un fichier et recueillir des données réelles}
Les valeurs dépendent de la collecte réelle. Vérifier : une catégorie par variable et par réponse ; somme des cellules égale au nombre de réponses exploitables ; dénominateur correct pour chaque fréquence ; arrondi annoncé. L'échantillon est une classe particulière, pas un échantillon représentatif de tous les élèves du pays. Les non-réponses doivent être distinguées d'une modalité observée.
\end{frame}
\end{document}
